\شروع‌مسئله[30] \متن‌سیاه{پیشنهاد الگوریتم یا داده‌ساختار مناسب}. در هر قسمت از این سؤال از شما خواسته می‌شود یک الگوریتم یا داده‌ساختار مناسب پیشنهاد دهید. این داده‌ساختار یا الگوریتم می‌تواند از مطالب گفته شده در درس، ترکیبی از آن‌ها و یا خلاقیتی باشد. علاوه بر \زیرخط{کلیت داده‌ساختار و الگوریتم}، \زیرخط{زمان اجرا} و \زیرخط{دلیل انتخاب خود} را نیز در ۳ الی ۵ خط به طور واضح توضیح دهید. در این سؤال نیاز به اثبات درستی الگوریتم و اثبات زمان اجرای الگوریتم نیست. همچنین در صورتی که زمان اجرای الگوریتم شما \زیرخط{زمان سرشکن یا زمان تصادفی} است، حتماً آن را بیان کنید.
% صورت اصلی مسئله را در این‌جا بنویسید.
\شروع{زیرمسائل}
\زیرمسئله[5]
یک آرایه از اعداد داریم که از ابتدا تا جایی از آن صعودی و پس از آن نزولی است. چه الگوریتمی را برای مرتب کردن این آرایه مناسب است؟ زمان اجرای آن چقدر است؟

\پاسخ{
	ابتدا با یک پیمایش عضو بیشینه‌ی آرایه را پیدا می‌کنیم، سپس اگر از اول آرایه تا این عضو را در یک آرایه در نظر بگیریم، و اعداد پس از این عضو را نیز به ترتیب برعکس در یک آرایه دیگر در نظر بگیریم، کافی است این ۲ قسمت را با هم ادغام کنیم‌(مانند مرتب‌سازی ادغامی)
	از آنجایی که ادغام از
	$O(n)$
	است، زمان اجرای الگوریتم ما نیز از 
	$O(n)$
	خواهد بود.
	
}
\نمره‌دهی{
	الگوریتم بالا نمره‌ی کامل را می‌گیرد و الگوریتم‌های دیگر با زمان اجرای بیش‌تر و راه‌حل‌های ناقص، متناسب با راه‌حل بخشی از نمره را گرفته‌اند.
	
}

\زیرمسئله[5]
یک آرایه از اعداد داریم که اندیس هر عدد با اندیسی که پس از مرتب شدن کل آرایه پیدا می‌کند، حداکثر ۱۰ تا اختلاف دارد. بهترین الگوریتم برای مرتب‌سازی این آرایه از بین الگوریتم‌های مرتب‌سازی که در کلاس مطرح شد کدام است؟ زمان اجرای آن چقدر است؟

\پاسخ{
	در صورتی که از مرتب سازی درجی استفاده کنیم، هر عضو زمانی که به آن می‌رسیم نیز حداکثر ۱۰ خانه با جایگاه اصلی خود فاصله دارد، پس در
	$O(1)$
	به جایگاه خود می‌رسد،
	در نتیجه در مجموع زمان اجرا از 
	$O(n)$
	خواهد بود.
	
}
\نمره‌دهی{
	اگر الگوریتم ارائه شده درست و از 
	$O(n)$
	باشد، نمره‌ی کامل را می‌گیرد، در غیر اینصورت متناسب با راه‌حل نمره می‌گیرد.
	
}

\زیرمسئله[8]
یک آرایه به نام $A$ و به طول $n$ از اعداد داریم که لزوماً مرتب نیست. یک عدد $x$ و یک مقدار صحیح و مثبت $k<n$ به ما داده شده است و می‌خواهیم $k$ تا بزرگترین عدد موجود در $A$ که کوچک‌تر از $x$ هستند را در صورت وجود پیدا کنیم. توجه کنید $x$ لزوماً در آرایه موجود نیست و ترتیب اعداد خروجی برای ما مهم نیست.

\پاسخ{
ابتدا با یک دور پیمایش آرایه اعدادی که بزرگ‌تر مساوی x هستند را حذف می‌کنیم سپس باید در آرایه باقی‌مانده k تا بزرگ‌ترین عدد را پیدا کنیم که این کار را هم در بهترین حالت می‌توان در زمان اجرای nlogk انجام داد. به این گونه که روی اعداد آرایه ثانویه پیمایش می‌کنیم و  هر عدد را به یک هرم کمینه اضافه می‌کنیم و اگر اندازه درخت بیشتر از k شده بود عضو کمینه را از آن حذف می‌کنیم. به این ترتیت در هر مرحله k تا عضو بیشینیه در هرم هستند و در نهایت می‌توانیم این اعداد را خروجی دهیم.
}

\نمره‌دهی{
	اگر راه‌حل گفته شده 
	$O(nlogk)$
	باشد، ۸ نمره تعلق می‌گیرد. اگر راه‌حل اشتباه \متن‌لاتین{counting sort} گفته شده باشد، ۲ نمره و اگر راه‌حل 
	$O(nlogn)$
	باشد هم ۱ نمره تعلق می‌گیرد.
	
}


%
%\زیرمسئله[8]
%یک مجموعه از اعداد متمایز داریم و به دنبال سه تا از این اعداد هستیم که جمع آن‌ها برابر با صفر شود. الگوریتمی با زمان اجرای $O(n^2)$ طراحی کنید که یک سه‌تایی با این خاصیت (در صورت وجود) پیدا کند. 
%به طراحی الگوریتمی با زمان اجرای $O(n^2\log n)$ نیمی از نمره داده خواهد شد. (توجه کنید شما مجاز به استفاده از درهم‌سازی نیستید)
%
%\پاسخ{
%عناصر را در $O(n\log n)$ مرتب می‌کنیم. همۀ عناصر را به عنوان عنصر اول سه‌تایی امتحان می‌کنیم. سپس یک اشاره‌گر مربوط به عنصر دوم به کوچکترین عنصر و یک اشاره‌گر مربوط به عنصر سوم به بزرگترین عنصر قرار می‌دهیم. در هر مرحله اگر جمع این سه عنصر بزرگ‌تر از صفر بود، اشاره‌گر عنصر سوم را حرکت می‌دهیم و اگر بزرگ‌تر از صفر بود، اشاره‌گر عنصر دوم را حرکت می‌دهیم. همچنین به تکراری نبودن عناصر سه‌تایی توجه می‌کنیم. زمان اجرای این الگوریتم برای یک عنصر اول $O(n)$ و در مجموع $O(n^2)$ می‌شود.
%	
%}

\زیرمسئله[6]
در این مسئله با یک داده ساختار جدید سر و کار داریم که مشابه هرم است ولی خاصیت «تقریباً کامل» بودن هرم را ندارد. در عوض مشابه درخت \متن‌لاتین{AVL}، می‌دانیم اختلاف ارتفاع زیردرخت فرزندان هر رأس حداکثر یک واحد با هم اختلاف دارد (مشابه درخت AVL ارتفاع رأس‌هایی که وجود ندارند را $-1$ درنظر می‌گیریم).
الگوریتمی برای درج یک عنصر جدید در این داده ساختار پیشنهاد دهید. 


\پاسخ{
	ابتدا مقدار جدید را به درخت اضافه می‌کنیم. سپس با استفاده از   buble up تا جای ممکن راس جدید را بالا می‌بریم. اگر در نهایت در مسیر طی شده همه راس‌ها ویژگی درخت AVL را داشتند یعنی درخت جدید هم ویژگی گفته شده را دارد. در غیر این صورت پایین‌ترین راسی که این ویژگی را دارد را در نظر می‌گیریم و مشابه rotate در درخت AVL ارتفاع درخت را یکی کم می‌کنیم.
	
}


\نمره‌دهی{
	اگر روش بیشتر از 
	$O(logn)$
	طول بکشد، نمره‌ای تعلق نمی‌گیرد. به ازای هرحالت کوچک که راه‌حل اشتباه کار کند یک نمره کم می‌شود.
	
}


\زیرمسئله[6]
$n$ 
شیء داریم که شیء $i$اُم دارای کلید $(a_i, b_i)$ است که $a_i$ها و $b_i$ اعداد صحیح و نامنفی کمتر از $n$ هستند. می‌خواهیم این اشیاء را مرتب کنیم. وقتی می‌خواهیم دو شیء را با هم مقایسه کنیم از تابع $f(a, b) = (1+\frac{a}{n})\cdot 2^b$ روی کلید هر کدام استفاده می‌کنیم و شیء‌ای کوچکتر است که خروجی $f$ روی کلید آن کوچکتر باشد. الگوریتمی برای مرتب‌سازی این اشیاء پیشنهاد دهید.


\پاسخ{
	ابتدا ثابت می‌کنیم که اگر دو تا 
	$b$
	نابرابر باشند مقدار
	$f$
	در آنی که
	$b$
	اش بزرگتر است بیشتر می‌باشد.
	
	$$
	b_1 < b_2 , a_1 < n, a_2 < n
	$$
	$$
	f(a_1, b_1) = (1 + \frac{a_1}{n})2^{b_1} < 2 \times 2^{b_1} \le 2^{b_2} \le (1 + \frac{a_2}{n})2^{b_2} = f(a_1, b_1)
	$$
	همچنین برای
	$b$
	های برابر باید طبق 
	$a$
	مقایسه شوند.
	
	برای حل می‌توانیم آن‌را به صورت عدد دو رقمی ببنیم در مبنای 
	$n$
	چرا که 
	$a, b < n$
	و از روش مرتب‌سازی مبنایی
	استفاده کنیم که
	$O(n)$
	به ما جوابی می‌دهد.
	
}

\نمره‌دهی{
	اگر به موضوع اولویت مقایسه
	$b$
	اشاره شده‌باشد ۳ نمره تعلق می‌گیرد.
	اگر به مرتب‌سازی مبنایی اشاره شده باشد هم ۳ نمره به شما تعلق می‌گیرد.
	
	راه‌حل‌هایی که الگوریتم با مرتبه زمانی بدتر از 
	$O(n)$
	پیشنهاد داده‌اند صفر در نظر گرفته‌می‌شوند.
	
	همچنین اشاره به اثبات اولویت مقایسه 
	$b$
	بخشی از نمره را دارد.
	
}


\پایان{زیرمسائل}